\section{Related Work}
\label{sec:related}

Our work draws from several distinct fields: authorship attribution, AI-generated text detection, and speech-to-text error analysis. We position our research as a bridge between these areas by introducing 'mode' as a unified latent variable.

{\bf Authorship Attribution.}
Traditional authorship attribution focuses on identifying individuals based on stylometric signatures \cite{kumarage2023neural}.
Recent advances have extended this to Large Language Models, identifying distinct writing styles and signatures inherent to different model architectures \cite{huang2024authorship}.
While these works focus on {\it who} wrote the text, our work focuses on {\it how} the text was produced (e.g., dictated vs. typed), regardless of the specific author.

{\bf AI-Generated Text Detection.}
The rapid adoption of LLMs has led to a surge in research on detecting AI-generated text.
State-of-the-art detectors like FAID \cite{ta2025faid} use contrastive learning to distinguish between human and AI content at fine-grained levels.
However, these detectors typically treat "human" as a monolithic category.
Our research decomposes the human category into 'Dictated' and 'Typed' modes, exploring the systematic artifacts that distinguish different human input methods.

{\bf ASR and OCR Artifacts.}
Forensic linguistics has long recognized that different input technologies leave systematic traces.
Speech-to-text (ASR) systems are known to produce phonetic errors and homophone substitutions that can derail downstream tasks like code understanding \cite{havare2026lost}.
Similarly, keyboard typing introduces orthographic errors based on physical key proximity.
We build on these observations by investigating whether these artifacts are sufficiently structured to be isolated as latent variables in modern embedding spaces.

{\bf Latent Variable Analysis.}
The conceptualization of language features as latent variables is gaining traction for reasoning optimization and model interpretability \cite{wu2026language}.
We extend this perspective to the origin of text, testing the hypothesis that 'mode' is a latent variable that modulates the representation of text in high-dimensional embedding spaces.
Unlike previous work that often requires supervised fine-tuning to isolate such variables, we focus on zero-shot detectability and inherent clustering in pre-trained models.
